\item[\textbf{Problem 13}]
What is the difference between using proxy server and the Tor?

\Solution

A proxy server functions as an intermediary between a client device and the
internet, forwarding requests such that the destination server observes the
proxy's IP address rather than the originating device's. This architecture,
however, introduces a single point of trust: the proxy operator retains
visibility into the client's original IP address, browsing destinations, and any
unencrypted traffic. While proxies offer low latency and are effective for
circumventing geo-restrictions or network-level access controls, they provide
limited anonymity and presuppose confidence in the operator's integrity.

Tor is designed with a fundamentally different threat model, prioritizing
trustless anonymity over performance. Rather than routing traffic through a
single intermediary, Tor applies multiple layers of encryption and distributes
the connection across at least three volunteer-operated relays — an entry,
middle, and exit node. Each relay decrypts only the layer necessary to identify
the next hop, ensuring that no individual node possesses both the client's
identity and the requested destination. This multi-hop architecture
substantially mitigates the risks of traffic analysis and network surveillance,
at the cost of increased latency inherent to the additional routing complexity.
